\begin{enumerate}
\item $a=-1$, $b=1$, $w(x)=1/\sqrt{1-x^2}$. 这给出 Chebyshev 多项式
\begin{equation}\label{eq:chebyshev-polynomial-cosine-formula}
T_n(x)=\dfrac1{2^{n-1}}\cos(n\arccos x).
\end{equation}

\item $a=-\infty$, $b=\infty$, $w(x)=\exp(-x^2)$. 这给出 Hermite 多项式
\begin{equation}\label{eq:hermite-polynomial-rodrigues-formula}
H_n(x)=(-1)^n\exp(x^2)\dfrac{\mathrm{d}^n}{\mathrm{d}x^n}\exp(-x^2).
\end{equation}

\item $a=0$, $b=\infty$, $w(x)=\exp(-x)$. 这给出 Laguerre 多项式
\begin{equation}\label{eq:laguerre-polynomial-rodrigues-formula}
L_n(x)=e^x\dfrac{\mathrm{d}^n}{\mathrm{d}x^n}(x^ne^{-x}).
\end{equation}
\end{enumerate}

还有其他各种具有专门名称的正交多项式, 例如 Jacobi 和 Gegenbauer 多项式. Chebyshev, Hermite 和 Laguerre 多项式的正交性留给读者验证; 见 Problem~\ref{ex:classical-polynomials-orthogonality-normalization}. 关于正交多项式已有许多已知事实, 例如系数公式, ``生成函数'', 正交多项式满足的微分方程, 递推关系以及与各种特殊函数的联系. 这里不讨论这些问题, 请参阅有关文献.

还需讨论以上介绍的多项式的完备性. 对 Chebyshev 多项式, 由 Weierstrass 逼近定理便显然可知, 因为一致收敛意味着 $L_w^2(-1,1)$ 中的收敛. 但对 Hermite 和 Laguerre 多项式, 处理的是无穷区间, 需要略有不同的论证. 先考察 Laguerre 多项式. 有恒等式
\begin{equation}\label{eq:laguerre-polynomial-generating-function}
\sum_{n=0}^\infty\dfrac{L_n(x)}{n!}t^n=g(x,t):=\dfrac1{1-t}\exp\left(-\dfrac{xt}{1-t}\right);
\end{equation}
见 Problem~\ref{ex:laguerre-generating-function}. 由于 Laguerre 多项式是由 \eqref{eq:laguerre-polynomial-generating-function} 右端的 Taylor 展开 ``生成'' 的, 这个表达式称为生成函数. 注意, 该 Taylor 级数的收敛半径为 1. 显式计算表明
\begin{equation}\label{eq:laguerre-generating-function-lp-two-remainder}
\int_0^\infty e^{-x}\left(g(x,t)-\sum_{n=0}^N\dfrac{L_n(x)}{n!}t^n\right)^2\,\mathrm{d}x=\dfrac1{1-t^2}-\sum_{n=0}^N t^{2n}
\end{equation}


(见 Problem~\ref{ex:laguerre-hermite-completeness-details}); 因而当 $|t|<1$ 时, \eqref{eq:laguerre-polynomial-generating-function} 中的级数也在 $L_w^2(0,\infty)$ 中收敛. 由此可得, 任意函数 $e^{-\alpha x}$, $\alpha>-1/2$, 都可在 $L_w^2(0,\infty)$ 中由 Laguerre 多项式逼近. 然而, 若 $f\in L_w^2$ 对每个 $\alpha$ 都与 $e^{-\alpha x}$ 正交, 则 $f$ 的 Laplace 变换为零, 从而 $f$ 为零. 因而指数函数 $e^{-\alpha x}$ 的线性组合在 $L_w^2(0,\infty)$ 中稠密.

对 Hermite 多项式, 有恒等式
\begin{equation}\label{eq:hermite-polynomial-generating-function}
e^{-t^2+2tx}=\sum_{n=0}^\infty\dfrac{H_n(x)}{n!}t^n
\end{equation}
与 \eqref{eq:laguerre-polynomial-generating-function} 对应, 并可应用类似论证. 此时, Taylor 级数的收敛半径为无穷, 需要证明函数 $e^{2tx}$, $t\in\mathbb{C}$, 的线性组合在 $L_w^2(-\infty,\infty)$ 中稠密. 首先注意 $\mathcal{D}(\mathbb{R})$ 稠密. $\mathcal{D}(\mathbb{R})$ 中的每个函数 $\phi$ 都可以由收敛的 Fourier 积分表示:
\begin{equation}\label{eq:hermite-completeness-fourier-inversion}
\phi(x)=\dfrac1{\sqrt{2\pi}}\int_{-\infty}^\infty e^{i\xi x}\widehat\phi(\xi)\,\mathrm{d}\xi,
\end{equation}
该积分可以由 Riemann 和逼近
\begin{equation}\label{eq:hermite-completeness-fourier-riemann-sums}
\phi(x)\sim\dfrac1{\sqrt{2\pi}}\sum_{k=1}^n e^{i\xi_kx}\widehat\phi(\xi_k)(\xi_k-\xi_{k-1}).
\end{equation}
容易看出, 这些离散和在 $L_w^2(-\infty,\infty)$ 中的收敛意义下收敛到该积分.
